\textbf{Phần 1. Trắc nghiệm lựa chọn}
\setcounterpageref{ex}{0}

\begin{ex} %[5D1N9]
Giá trị của chữ số $6$ trong số $960\,730$ là
\choice
{\True $60\,000$}
{$60$}
{$600$}
{$6\,000$}
\loigiai{
Trong số $960\,730$, chữ số $6$ ở hàng chục nghìn nên có giá trị là $60\,000$.\\
\textbf{Đáp án của câu hỏi là A}
}
\end{ex}

\begin{ex} %[5D1N9]
Số lớn nhất trong các số $109\,989; 105\,789; 110\,200; 99\,000$ là
\choice
{$109\,989$}
{$105\,789$}
{\True $110\,200$}
{$99\,000$}
\loigiai{
So sánh các số: $99\,000 < 105\,789 < 109\,989 < 110\,200$. Vậy số lớn nhất là $110\,200$.\\
\textbf{Đáp án của câu hỏi là C}
}
\end{ex}

\begin{ex} %[5D1H9]
Phân số bé nhất trong các phân số $\dfrac{23}{36}; \dfrac{11}{12}; \dfrac{7}{9}; \dfrac{2}{3}$ là
\choice
{$\dfrac{11}{12}$}
{$\dfrac{7}{9}$}
{$\dfrac{2}{3}$}
{\True $\dfrac{23}{36}$}
\loigiai{
Quy đồng mẫu số chung là $36$:\\
$\dfrac{11}{12} = \dfrac{33}{36}$; $\dfrac{7}{9} = \dfrac{28}{36}$; $\dfrac{2}{3} = \dfrac{24}{36}$; $\dfrac{23}{36}$.\\
Vì $\dfrac{23}{36} < \dfrac{24}{36} < \dfrac{28}{36} < \dfrac{33}{36}$ nên phân số bé nhất là $\dfrac{23}{36}$.\\
\textbf{Đáp án của câu hỏi là D}
}
\end{ex}

\begin{ex} %[5D1H9]
Phân số $\dfrac{143}{100}$ viết dưới dạng hỗn số là
\choice
{\True $1\dfrac{43}{100}$}
{$14\dfrac{3}{100}$}
{$140\dfrac{3}{100}$}
{$3\dfrac{14}{100}$}
\loigiai{
Ta có: $143 : 100 = 1$ dư $43$. Do đó $\dfrac{143}{100} = 1\dfrac{43}{100}$.\\
\textbf{Đáp án của câu hỏi là A}
}
\end{ex}

\begin{ex} %[5D1H9]
Giá trị của biểu thức $35\,700 : 50 + 68 \times 46$ là
\choice
{$3\,128$}
{$3\,742$}
{\True $3\,842$}
{$4\,128$}
\loigiai{
Ta có: $35\,700 : 50 = 714$ và $68 \times 46 = 3\,128$.\\
Tổng là: $714 + 3\,128 = 3\,842$.\\
\textbf{Đáp án của câu hỏi là C}
}
\end{ex}

\begin{ex} %[5D1H9]
Giá trị của biểu thức $\dfrac{16}{9} \times \left(6 : \dfrac{32}{9}\right) - \dfrac{7}{5}$ là
\choice
{$3$}
{\True $\dfrac{8}{5}$}
{$\dfrac{7}{5}$}
{$\dfrac{1}{5}$}
\loigiai{
Ta có: $6 : \dfrac{32}{9} = 6 \times \dfrac{9}{32} = \dfrac{27}{16}$.\\
Khi đó: $\dfrac{16}{9} \times \dfrac{27}{16} - \dfrac{7}{5} = 3 - \dfrac{7}{5} = \dfrac{8}{5}$.\\
\textbf{Đáp án của câu hỏi là B}
}
\end{ex}

\begin{ex} %[5D1H9]
Viết số đo $52\text{ cm}$ dưới dạng phân số thập phân với đơn vị mét là
\choice
{\True $\dfrac{52}{100}\text{ m}$}
{$\dfrac{52}{10}\text{ m}$}
{$\dfrac{52}{1\,000}\text{ m}$}
{$\dfrac{52}{100}\text{ cm}$}
\loigiai{
Vì $1\text{ m} = 100\text{ cm}$ nên $52\text{ cm} = \dfrac{52}{100}\text{ m}$.\\
\textbf{Đáp án của câu hỏi là A}
}
\end{ex}

\begin{ex} %[5D1H9]
Viết số đo $2\text{ dm } 7\text{ cm}$ dưới dạng hỗn số với đơn vị đề-xi-mét là
\choice
{$7\dfrac{2}{10}\text{ dm}$}
{$27\dfrac{1}{10}\text{ dm}$}
{$2\dfrac{7}{100}\text{ dm}$}
{\True $2\dfrac{7}{10}\text{ dm}$}
\loigiai{
Ta có: $2\text{ dm } 7\text{ cm} = 2\text{ dm} + \dfrac{7}{10}\text{ dm} = 2\dfrac{7}{10}\text{ dm}$.\\
\textbf{Đáp án của câu hỏi là D}
}
\end{ex}

\begin{ex} %[5D1H9]
Tính nhanh biểu thức $524 \times 63 + 524 \times 37 - 2\,400$ ta được
\choice
{$5\,000$}
{\True $50\,000$}
{$52\,400$}
{$500\,000$}
\loigiai{
$524 \times (63 + 37) - 2\,400 = 524 \times 100 - 2\,400 = 52\,400 - 2\,400 = 50\,000$.\\
\textbf{Đáp án của câu hỏi là B}
}
\end{ex}

\begin{ex} %[5D1V9]
Một bút bi có giá $4\,500$ đồng, một quyển vở có giá $7\,000$ đồng. Nam mua $2$ bút bi và $7$ quyển vở rồi đưa cho cô bán hàng tờ tiền $100\,000$ đồng. Cô bán hàng trả lại Nam số tiền là
\choice
{$58\,000$ đồng}
{$48\,000$ đồng}
{\True $42\,000$ đồng}
{$52\,000$ đồng}
\loigiai{
Số tiền Nam mua bút và vở là: $2 \times 4\,500 + 7 \times 7\,000 = 9\,000 + 49\,000 = 58\,000$ đồng.\\
Số tiền cô bán hàng trả lại là: $100\,000 - 58\,000 = 42\,000$ đồng.\\
\textbf{Đáp án của câu hỏi là C}
}
\end{ex}

\begin{ex} %[5D1V9]
Một đoàn xe gồm $6$ xe ô tô loại nhỏ chở $35$ học sinh/xe và $9$ xe loại lớn chở $40$ học sinh/xe. Trung bình mỗi xe chở bao nhiêu học sinh?
\choice
{$37$ học sinh}
{$39$ học sinh}
{$37{,}5$ học sinh}
{\True $38$ học sinh}
\loigiai{
Tổng số học sinh là: $6 \times 35 + 9 \times 40 = 210 + 360 = 570$ học sinh.\\
Tổng số xe là: $6 + 9 = 15$ xe.\\
Trung bình mỗi xe chở: $570 : 15 = 38$ học sinh.\\
\textbf{Đáp án của câu hỏi là D}
}
\end{ex}

\begin{ex} %[5D1V9]
Cô Ba mang $120$ quả trứng ra chợ bán. Lần thứ nhất cô bán được $\dfrac{1}{8}$ số trứng đó. Lần thứ hai cô bán được $\dfrac{2}{7}$ số trứng còn lại sau lần bán thứ nhất. Hỏi sau cả hai lần, cô Ba đã bán được tất cả bao nhiêu quả trứng?
\choice
{$30$ quả}
{\True $45$ quả}
{$75$ quả}
{$40$ quả}
\loigiai{
Lần thứ nhất bán được: $120 \times \dfrac{1}{8} = 15$ (quả trứng).\\
Số trứng còn lại là: $120 - 15 = 105$ (quả trứng).\\
Lần thứ hai bán được: $105 \times \dfrac{2}{7} = 30$ (quả trứng).\\
Cả hai lần bán được: $15 + 30 = 45$ (quả trứng).\\
\textbf{Đáp án của câu hỏi là B}
}
\end{ex}

\textbf{Phần 2. Trắc nghiệm đúng sai}
\setcounterpageref{ex}{0}

\begin{ex} %[5D1H9]
Cho các biểu thức: $A = \dfrac{7}{9} \times \dfrac{6}{5} + \dfrac{6}{5} \times \dfrac{2}{9} + \dfrac{4}{5}$ và $B = 524 \times 63 + 524 \times 37 - 2\,400$. Xét tính đúng sai của các khẳng định sau:
\choiceTF
{\True Biểu thức $A$ có giá trị bằng $2$.}
{Biểu thức $B$ có giá trị bằng $5\,000$.}
{\True Trong biểu thức $A$, áp dụng tính chất phân phối ta có $\dfrac{6}{5} \times \left(\dfrac{7}{9} + \dfrac{2}{9}\right) = \dfrac{6}{5}$.}
{\True Giá trị của biểu thức $B$ gấp $25\,000$ lần giá trị của biểu thức $A$.}
\loigiai{
\begin{itemize}
    \item $A = \dfrac{6}{5} \times \left(\dfrac{7}{9} + \dfrac{2}{9}\right) + \dfrac{4}{5} = \dfrac{6}{5} \times 1 + \dfrac{4}{5} = \dfrac{10}{5} = 2$.
    \item $B = 524 \times 100 - 2\,400 = 52\,400 - 2\,400 = 50\,000 \neq 5\,000$.
    \item Ta có $\dfrac{6}{5} \times 1 = \dfrac{6}{5}$.
    \item Tỉ số $B : A = 50\,000 : 2 = 25\,000$.
\end{itemize}
\begin{itemize}
    \item a) \textbf{Đúng}.
    \item b) \textbf{Sai}: $B = 50\,000$.
    \item c) \textbf{Đúng}.
    \item d) \textbf{Đúng}.
\end{itemize}
}
\end{ex}

\begin{ex} %[5H1V9]
Sân trường hình chữ nhật có chu vi là $84\text{ m}$, chiều dài hơn chiều rộng là $8\text{ m}$. Xét tính đúng sai của các khẳng định sau:
\choiceTF
{\True Nửa chu vi sân trường là $42\text{ m}$.}
{Chiều rộng sân trường là $25\text{ m}$.}
{\True Chiều dài sân trường là $25\text{ m}$.}
{\True Diện tích của sân trường đó là $425\text{ m}^2$.}
\loigiai{
\begin{itemize}
    \item Nửa chu vi là: $84 : 2 = 42\text{ (m)}$.
    \item Chiều dài sân trường là: $(42 + 8) : 2 = 25\text{ (m)}$.
    \item Chiều rộng sân trường là: $25 - 8 = 17\text{ (m)} \neq 25\text{ m}$.
    \item Diện tích sân trường là: $25 \times 17 = 425\text{ (m}^2\text{)}$.
\end{itemize}
\begin{itemize}
    \item a) \textbf{Đúng}.
    \item b) \textbf{Sai}: Chiều rộng là $17\text{ m}$.
    \item c) \textbf{Đúng}.
    \item d) \textbf{Đúng}.
\end{itemize}
}
\end{ex}

\textbf{Phần 3. Trả lời ngắn}
\setcounterpageref{ex}{0}

\begin{ex} %[5D1H9]
Tính giá trị của biểu thức sau bằng cách thuận tiện nhất:
\[
M = \dfrac{7}{9} \times \dfrac{6}{5} + \dfrac{6}{5} \times \dfrac{2}{9} + \dfrac{4}{5}.
\]
\shortans{2}
\loigiai{
\textbf{Đáp số:} $2$\\
$M = \dfrac{6}{5} \times \left(\dfrac{7}{9} + \dfrac{2}{9}\right) + \dfrac{4}{5} = \dfrac{6}{5} \times 1 + \dfrac{4}{5} = \dfrac{10}{5} = 2$.
}
\end{ex}

\begin{ex} %[5D1H9]
Ước lượng kết quả phép tính $4\,180\,300 + 3\,990\,700$ đến hàng triệu. Kết quả ước lượng là bao nhiêu triệu?
\shortans{8}
\loigiai{
\textbf{Đáp số:} $8$\\
$4\,180\,300 \approx 4\,000\,000$; $3\,990\,700 \approx 4\,000\,000 \Rightarrow 4 + 4 = 8$ (triệu).
}
\end{ex}

\begin{ex} %[5D1H9]
Trạng Trình Nguyễn Bỉnh Khiêm sinh năm 1491. Hỏi năm kỉ niệm 600 năm năm sinh của Trạng Trình là năm nào?
\shortans{2091}
\loigiai{
\textbf{Đáp số:} $2091$\\
Năm kỉ niệm là: $1491 + 600 = 2091$.
}
\end{ex}

\begin{ex} %[5H1V9]
Sân trường hình chữ nhật có chu vi $84\text{ m}$, chiều dài hơn chiều rộng $8\text{ m}$. Tính diện tích của sân trường đó theo mét vuông.
\shortans{425}
\loigiai{
\textbf{Đáp số:} $425$\\
Nửa chu vi: $84 : 2 = 42\text{ m}$.\\
Chiều dài: $(42 + 8) : 2 = 25\text{ m}$, Chiều rộng: $42 - 25 = 17\text{ m}$.\\
Diện tích: $25 \times 17 = 425\text{ m}^2$.
}
\end{ex}

\begin{ex} %[5D1V9]
Cô Ba mang $120$ quả trứng gà ra chợ bán. Lần 1 bán được $\dfrac{1}{8}$ số trứng. Lần 2 bán được $\dfrac{2}{7}$ số trứng còn lại sau lần 1. Hỏi sau hai lần bán, cô Ba còn lại bao nhiêu quả trứng?
\shortans{75}
\loigiai{
\textbf{Đáp số:} $75$\\
Lần 1 bán: $120 \times \dfrac{1}{8} = 15$ quả.\\
Còn lại sau lần 1: $120 - 15 = 105$ quả.\\
Lần 2 bán: $105 \times \dfrac{2}{7} = 30$ quả.\\
Số trứng còn lại sau cả hai lần: $105 - 30 = 75$ quả.
}
\end{ex}

\begin{ex} %[5D1V9]
Một đoàn xe gồm $6$ xe ô tô chở $35$ học sinh/xe và $9$ xe chở $40$ học sinh/xe. Hỏi trung bình mỗi xe chở bao nhiêu học sinh?
\shortans{38}
\loigiai{
\textbf{Đáp số:} $38$\\
Tổng học sinh: $6 \times 35 + 9 \times 40 = 210 + 360 = 570$ em.\\
Trung bình mỗi xe: $570 : 15 = 38$ em.
}
\end{ex}
